% =============================================================
% MAJOR-REVISION RESPONSE-TO-REVIEWERS TEMPLATE
% =============================================================
% Generic LaTeX scaffold for a major-revision response document,
% suitable for OOPSLA/ICSE/FSE/ASE/ISSTA-style submissions.
% Adapt to your paper:
%   1. Replace <TODO ...> placeholders with the real content.
%   2. Keep one \section per reviewer (Meta + A + B + C + ...).
%   3. Order items inside each section as the reviewer wrote them.
%   4. For every promised change, immediately follow with the revised
%      paper prose in a \shadedquotation block.
%   5. Use \hyperref[resp:<label>]{...} to cross-link to other items
%      whenever the same concern appears in multiple places.
%
% Conventions (style varies per venue/paper — adapt as needed):
%   \pointRaised{<ReviewerID>}{<Quoted reviewer comment>}
%   \reply{\responsed{<Your response, 2–4 sentences.>}}
%   \revised{In Section~X, we added:}
%   \begin{shadedquotation} <new paper text> \end{shadedquotation}
%
% =============================================================

\documentclass[12pt]{article}
\usepackage{url}
\usepackage{vmargin}
\setpapersize{USletter}
\setmarginsrb{1in}{1in}{1in}{0.5in}{0pt}{0mm}{10pt}{0.5in}
\usepackage{charter}
\usepackage{color}
\usepackage{graphicx}
\usepackage{booktabs}
\usepackage{xspace}
\usepackage{amsmath,amssymb,amsfonts}
\usepackage[many]{tcolorbox}
\usepackage{textcomp}
\usepackage{xcolor}
\usepackage{enumitem}
\usepackage{multirow}
\usepackage{hyperref}
\usepackage{pifont}
\usepackage{cleveref}
\usepackage{framed}
\usepackage{quoting}
\usepackage[x11names]{xcolor}
\usepackage{caption}

\setlength{\parskip}{1ex plus1ex minus1ex}
\setlength{\parindent}{0pt}

% --- Core macros (lift these as-is) ---------------------------
\newcommand{\pointRaised}[2]{\smallskip
{{\fontseries{b}\selectfont #1}: \textbf{#2}}\newline}
\newcommand{\simplePointRaised}[1]{\bigskip \hrule\textsl{#1} }
\newcommand{\reply}[1]{\textbf{{Response}}:\ {#1} \smallskip }
\newcommand{\rev}[1]{\textcolor{black}{#1}}
\newcommand{\revise}[1]{\textcolor{red}{#1}}
\newcommand{\revised}[1]{\textcolor{blue}{#1}}
\renewcommand{\rev}[1]{{\color{black}#1}}
\renewcommand{\revise}[1]{{\color{red}#1}}
\renewcommand{\revised}[1]{{\color{blue}#1}}
\newcommand{\responsed}[1]{{\color{black}#1}}

% Your tool/project name macro — rename or remove
\newcommand{\tool}{\textit{<TOOL_NAME>}\xspace}
\newcommand{\name}{\tool}

\newcommand{\myfig}{Fig.\xspace}
\newcommand{\mysec}{\S}
\newcommand{\fixme}[1]{\textcolor{red}{#1}}

% "How to read" callout box
\DeclareRobustCommand{\mybox}[2][gray!20]{%
\begin{tcolorbox}[
        breakable,
        left=0pt, right=0pt, top=0pt, bottom=0pt,
        colback=#1, colframe=#1,
        width=\linewidth,
        enlarge left by=0mm,
        boxsep=5pt,
        arc=0pt, outer arc=0pt,
        ]
        #2
\end{tcolorbox}
}

% "Revised paper text" environment — shaded box around quoted new prose
\colorlet{shadecolor}{LavenderBlush2}
\newenvironment{shadedquotation}
 {\begin{shaded*}\quoting[leftmargin=0pt, vskip=0pt]}
 {\endquoting\end{shaded*}}

% =============================================================

\title{Response to Reviewers' Comments\\[0.3em]
\large <TODO: Paper title here>
}

\begin{document}
\maketitle

\textit{We sincerely thank the reviewers and the program committee for their
careful reading of our paper and their constructive comments. We have
thoroughly revised the manuscript to address all points raised. Below, we
provide our point-by-point response in the following format:}

\tableofcontents
\bigskip

% --- "How to read this rebuttal" callout ---------------------
\mybox[red!8]{
{\Large\textcolor{red!70!black}{\textbf{How to read this rebuttal}}}

\begin{quote}
\pointRaised{Reviewer-A}{Comments or Questions from Reviewer A}
\end{quote}

\begin{quote}
\reply{\responsed{Response to Reviewer A's comments and questions.}}
\end{quote}

\begin{quote}
\revised{In Section~x, we added:}
\begin{shadedquotation}
\responsed{Revised manuscript text appears in this shaded block.}
\end{shadedquotation}
\end{quote}
}

%% ============================================================
%% META-REVIEW (REQUIRED CHANGES)
%% ============================================================
\section{Response to Meta-Review}
\label{sec:meta}

% --- Meta-R1 -------------------------------------------------
\phantomsection
\label{resp:meta-r1}
\pointRaised{Meta-R1}{<TODO: Quote the first meta-review item verbatim.>}

\reply{\responsed{%
<TODO: 2–4 sentences. Acknowledge the point, answer it directly, name the
section that will change.>
}}

\revised{In Section~<X.Y>, we added:}
\begin{shadedquotation}
\responsed{<TODO: Paste the new paper prose verbatim. This should match
exactly what now appears in the revised PDF (in blue \revision{...}).>}
\end{shadedquotation}

% Optionally include a table/figure that you added to the paper
% \begin{table}[h]
% \centering
% \caption{<Caption matching the new table in the paper.>}
% \label{tab:<new-table>}
% \begin{tabular}{...}
%   ...
% \end{tabular}
% \end{table}

% --- Meta-R2 -------------------------------------------------
\bigskip
\phantomsection
\label{resp:meta-r2}
\pointRaised{Meta-R2}{<TODO: Quote the second meta-review item verbatim.>}

\reply{\responsed{%
<TODO: Response. Use \\\\ between paragraphs inside the reply when needed.
Use \begin{itemize}...\end{itemize} for enumerated sub-arguments — wrap each
item body in \responsed{...} so the colour stays consistent.>
\\\\
\textbf{(1) <Sub-point header>.} <Sub-point body.>
\\\\
\textbf{(2) <Sub-point header>.} <Sub-point body.>
\\\\
\textbf{(3) <Sub-point header>.} <Sub-point body.>
\begin{itemize}
\item \responsed{\textbf{<Bullet header>:} <Bullet body, including numeric
results when relevant.>}
\item \responsed{\textbf{<Bullet header>:} <Bullet body.>}
\end{itemize}
\responsed{<Optional closing paragraph that ties the bullets together and
references prior work or sibling responses.>}
}}

\revised{In Section~<X.Y>, we added:}
\begin{shadedquotation}
\responsed{<TODO: New paper prose.>}
\end{shadedquotation}

% --- Cross-link pattern --------------------------------------
% When the same concern is raised in two places, cross-reference instead of
% repeating yourself:
%   See also \hyperref[resp:comment-c1]{Comment~C1} for related discussion.

% --- Repeat for Meta-R3, Meta-R4, ... ------------------------


%% ============================================================
%% REVIEWER A (Overall merit: <score, short label>)
%% ============================================================
\section{Response to Reviewer A}
\label{sec:reviewer-a}

% --- Reviewer's explicit questions first ---------------------
\bigskip
\phantomsection
\label{resp:question-aq1}
\pointRaised{Question A-Q1}{<TODO: Quote the question verbatim.>}

\reply{\responsed{%
<TODO: Direct answer, 2–4 sentences. If you have a concrete number that
answers the question, lead with it.>
}}

\revised{In <Section / RQ>, we added:}
\begin{shadedquotation}
\responsed{<TODO: New paper prose.>}
\end{shadedquotation}

\bigskip
\phantomsection
\label{resp:question-aq2}
\pointRaised{Question A-Q2}{<TODO: Quote question.>}

\reply{\responsed{%
<TODO: Response.>
}}

\revised{In Section~<X.Y>, we added:}
\begin{shadedquotation}
\responsed{<TODO: New paper prose.>}
\end{shadedquotation}

% --- Then the reviewer's detailed comments -------------------
\bigskip
\phantomsection
\label{resp:comment-a3}
\pointRaised{Comment A3}{<TODO: Quote the reviewer's longer comment
verbatim. Even if only part is actionable, keep the whole comment — the
surrounding framing shows you understood the point.>}

\reply{\responsed{%
<TODO: Response. Acknowledge the legitimate part of the critique even when
you don't fully concede. If you concede, narrow the claim; don't hedge it
with a three-word qualifier.>
\\\\
<TODO: Optional second paragraph with numeric evidence or a worked example.
Cross-link to other items if relevant:>
See also \hyperref[resp:question-aq1]{Question~A-Q1} for the quantitative
breakdown.
}}

\revised{In Section~<X.Y>, we added:}
\begin{shadedquotation}
\responsed{<TODO: New paper prose, possibly multi-paragraph. If the new
content is large, split into multiple \begin{shadedquotation}...
\end{shadedquotation} blocks anchored to different sections of the paper.>}
\end{shadedquotation}

% --- Repeat for Comment A4, A5, ... --------------------------


%% ============================================================
%% REVIEWER B (Overall merit: <score>)
%% ============================================================
\section{Response to Reviewer B}
\label{sec:reviewer-b}

\bigskip
\phantomsection
\label{resp:question-bq1}
\pointRaised{Question B-Q1}{<TODO: Quote.>}

\reply{\responsed{%
<TODO: Response.>
}}

\revised{In Section~<X.Y>, we added:}
\begin{shadedquotation}
\responsed{<TODO: New paper prose.>}
\end{shadedquotation}

% --- Repeat ---


%% ============================================================
%% REVIEWER C (Overall merit: <score>)
%% ============================================================
\section{Response to Reviewer C}
\label{sec:reviewer-c}

\bigskip
\phantomsection
\label{resp:question-cq1}
\pointRaised{Question C-Q1}{<TODO: Quote.>}

\reply{\responsed{%
<TODO: Response.>
}}

\revised{In Section~<X.Y>, we added:}
\begin{shadedquotation}
\responsed{<TODO: New paper prose.>}
\end{shadedquotation}

% --- Repeat ---


%% ============================================================
%% BIBLIOGRAPHY (only if response cites things)
%% ============================================================
\bibliographystyle{ACM-Reference-Format}
\bibliography{response}

\end{document}
